%! Unit 2: Describing One Variable — Summative Assessment — Unit Test (blank) (Answer Key)
% The GRADED form. It is never published into a packet (tests.mk drops only the practice test).
% A parallel form of practice_test: same file structure, parts, guards and answer-slot order;
% every number, context, vocabulary letter and correct MC letter differs.
%
% THE KEY IS GENERATED, NEVER HAND-EDITED:
%   python3 templates/lesson/mkkey.py unit02/tests/actual_test/main.tex \
%       unit02/test_keys/actual_test_key/main.tex unit02/test_keys/actual_test_key/answers.json
% Re-run it after every edit to this file. Part B carries an Answer: slot on every item (the
% quiz-01 model) so the letter lives in the generated key; unit 2 has no unit_cover_key page to
% park the letters on.
\documentclass[10pt]{article}
\usepackage{saar-article}
\usepackage{saar-key}   % pulls in -boxes; do NOT also load -boxes

% Work blocks on a test need handwriting room, not typeset spacing. This moves
% the blank and the key together, so it cannot open a page mismatch.
\setlength{\workrowsep}{5pt}

% Part-divider strip
\newcommand{\parthead}[1]{%
  \vspace{6pt}\noindent
  \begin{headlinebox}{cerulean}{\color{white}\bfseries #1}\end{headlinebox}%
  \vspace{4pt}}

% One horizontal boxplot on a tikz x-axis: \drawbox{y}{min}{Q1}{med}{Q3}{max}{label}
\newcommand{\drawbox}[7]{%
  \draw[charcoal, thick] (#2,#1) -- (#3,#1);
  \draw[charcoal, thick] (#5,#1) -- (#6,#1);
  \draw[charcoal, thick] (#2,#1-0.15) -- (#2,#1+0.15);
  \draw[charcoal, thick] (#6,#1-0.15) -- (#6,#1+0.15);
  \filldraw[fill=frost, draw=cerulean, thick] (#3,#1-0.3) rectangle (#5,#1+0.3);
  \draw[cerulean, very thick] (#4,#1-0.3) -- (#4,#1+0.3);
  \node[left, font=\small\bfseries] at (#2-0.2,#1) {#7};}

\begin{document}

\pageheader{Unit 2: Describing One Variable}{Unit Test --- Answer Key}
\namedateperiod

\vspace{0.10in}
\begin{remindbox}
\small
\textbf{Instructions:} Show all work. Answer every interpretation in context, with units.
No notes unless your teacher says so.
\textbf{25 questions, 100 points:} Part A $10$, Part B $12$, Part C $48$, Part D $30$.
\end{remindbox}

% ======================================================================
\parthead{Part A --- Vocabulary \hfill 10 questions, 1 point each}

\small
\textbf{Set 1 --- Displays and Shape.} Write the letter of the matching description
next to each term.

\vspace{2pt}
\begin{minipage}[t]{0.36\linewidth}
\begin{enumerate}[leftmargin=1.9em, itemsep=3.5pt, topsep=1pt]
  \item \ans{D} Distribution
  \item \ans{C} Stemplot
  \item \ans{A} Histogram
  \item \ans{E} Skewed right
  \item \ans{B} Skewed left
\end{enumerate}
\end{minipage}\hfill
\begin{minipage}[t]{0.61\linewidth}
\begin{enumerate}[label=\Alph*., leftmargin=1.9em, itemsep=3.5pt, topsep=1pt]
  \item Bars over equal-width bins; each bar's height is a count.
  \item The long tail stretches toward the smaller values.
  \item Each value split into a stem and a leaf, read with a key.
  \item What values a variable takes and how often it takes each one.
  \item The long tail stretches toward the larger values.
\end{enumerate}
\end{minipage}

\vspace{6pt}
\textbf{Set 2 --- Measures and Outliers.} Same directions.

\vspace{2pt}
\begin{minipage}[t]{0.36\linewidth}
\begin{enumerate}[leftmargin=1.9em, itemsep=3.5pt, topsep=1pt, start=6]
  \item \ans{B} Median
  \item \ans{A} Interquartile range
  \item \ans{E} Standard deviation
  \item \ans{D} Outlier
  \item \ans{C} Resistant
\end{enumerate}
\end{minipage}\hfill
\begin{minipage}[t]{0.61\linewidth}
\begin{enumerate}[label=\Alph*., leftmargin=1.9em, itemsep=3.5pt, topsep=1pt]
  \item $Q_3 - Q_1$: the spread of the middle half of the data.
  \item The middle value of the ordered data.
  \item Barely moved by one extreme value.
  \item A value below $Q_1 - 1.5 \times \text{IQR}$ or above
        $Q_3 + 1.5 \times \text{IQR}$.
  \item About the typical distance of the values from their mean.
\end{enumerate}
\end{minipage}
\normalsize

% ======================================================================
\boxguard[8]
\parthead{Part B --- Multiple Choice \hfill 6 questions, 2 points each}

\small
\begin{enumerate}[leftmargin=2.2em, itemsep=5pt, topsep=2pt, start=11]

\boxguard[9]
\item A histogram of $80$ Harlow students' scores on an easy review quiz has its tallest bars
between $90$ and $100$ points, then shorter and shorter bars down to $40$ points. The
distribution is
\begin{enumerate}[label=(\Alph*), leftmargin=1.8em, itemsep=1pt, topsep=1pt]
  \item skewed right, because the tallest bars are on the right.
  \item symmetric, because it has only one peak.
  \item skewed right, because most students scored high.
  \item skewed left, because the long tail stretches toward the smaller values.
\end{enumerate}
\hfill \textbf{Answer:} \ans{D}

\boxguard[9]
\item Salaries at a company are strongly skewed right by a few very large executive salaries.
The best measure of a typical salary is
\begin{enumerate}[label=(\Alph*), leftmargin=1.8em, itemsep=1pt, topsep=1pt]
  \item the mean, because it uses every salary.
  \item the range, because it shows the largest salary.
  \item the median, because the few large salaries barely move it.
  \item the standard deviation, because it measures typical distance.
\end{enumerate}
\hfill \textbf{Answer:} \ans{C}

\boxguard[9]
\item Relay team R's $400$-meter times run from $50$ to $54$ seconds with $s = 1.2$. Team S's
run from $58$ to $80$ seconds with $s = 7$. Which statement is true?
\begin{enumerate}[label=(\Alph*), leftmargin=1.8em, itemsep=1pt, topsep=1pt]
  \item Team S's times typically sit farther from their own mean than Team R's.
  \item Team R is faster, so its standard deviation must be smaller.
  \item The two standard deviations cannot be compared.
  \item Team S's times are larger, and that is why its $s$ is larger.
\end{enumerate}
\hfill \textbf{Answer:} \ans{A}

\boxguard[9]
\item Every value in a data set is decreased by $3$. What happens to the mean and the standard
deviation?
\begin{enumerate}[label=(\Alph*), leftmargin=1.8em, itemsep=1pt, topsep=1pt]
  \item Both decrease by $3$.
  \item The mean stays the same; the standard deviation decreases by $3$.
  \item Neither one changes.
  \item The mean decreases by $3$; the standard deviation does not change.
\end{enumerate}
\hfill \textbf{Answer:} \ans{D}

\boxguard[9]
\item One very small outlier is removed from a data set of $17$ values. Which summary
changes the most?
\begin{enumerate}[label=(\Alph*), leftmargin=1.8em, itemsep=1pt, topsep=1pt]
  \item the median
  \item the mean
  \item the third quartile
  \item the interquartile range
\end{enumerate}
\hfill \textbf{Answer:} \ans{B}

\boxguard[9]
\item Two clinics each have a median wait of $12$ minutes. Clinic M's IQR is $9$ minutes
and Clinic N's is $3$ minutes. Which is a correct comparison?
\begin{enumerate}[label=(\Alph*), leftmargin=1.8em, itemsep=1pt, topsep=1pt]
  \item Clinic M is more consistent, because its IQR is larger.
  \item The clinics perform the same, because their medians are equal.
  \item Both typically take about $12$ minutes, but Clinic N's waits are more consistent.
  \item Clinic N is faster, because its IQR is smaller.
\end{enumerate}
\hfill \textbf{Answer:} \ans{C}

\end{enumerate}
\normalsize

% ======================================================================
\boxguard[8]
\parthead{Part C --- Short Answer \& Computation \hfill 6 questions, 8 points each}

\small
\begin{enumerate}[leftmargin=2.2em, itemsep=9pt, topsep=2pt, start=17]

% ---------------------------------------------------------------- C1 (2.1)
\item Eighteen \textbf{Harlow High School} students recorded the hours they volunteered this
semester.

\vspace{2pt}
\begin{minipage}[t]{0.30\linewidth}
\vspace{0pt}
\renewcommand{\arraystretch}{1.1}
\begin{tabular}{r|l}
0 & 5\;7 \\
1 & 1\;2\;3\;3\;6\;7 \\
2 & 0\;1\;2\;2\;4\;5\;8 \\
3 & 1\;3 \\
4 & 6 \\
\end{tabular}

\vspace{3pt}
\textbf{Key:} $1\,|\,3$ means $13$ hours
\end{minipage}\hfill
\begin{minipage}[t]{0.67\linewidth}
\vspace{0pt}
(a) What does the leaf $6$ on stem $4$ stand for?

\vspace{2pt}
\ans{One student who volunteered 46 hours}

\vspace{4pt}
(b) Median: \ans{$20.5$ hours} \quad Range: \ans{$41$ hours}

\vspace{4pt}
(c) Describe the distribution --- shape, unusual features, center, and spread --- in
context.
\end{minipage}

\par\ansline{Roughly symmetric with a slight right tail, a gap from 33 to 46 hours, and a possible outlier at 46.}
\par\ansline{The median is 20.5 hours, and the hours run from 5 to 46 (range 41 hours).}

% ---------------------------------------------------------------- C3 (2.2)
\boxguard[22]
\item Nine members of the Harlow chess club won these numbers of games this season, in
order: $2$, $4$, $4$, $5$, $7$, $10$, $11$, $13$, $16$.

\vspace{2pt}
(a) Find the mean.

\begin{work}
  \bar{x} &= \dfrac{2 + 4 + 4 + 5 + 7 + 10 + 11 + 13 + 16}{9} = \dfrac{72}{9} = 8 \text{ games}
\end{work}

(b) Median: \ans{$7$ games} \quad Mode: \ans{$4$ games}

\vspace{3pt}
(c) Find $Q_1$ and $Q_3$. (Leave the median out of both halves.)

\begin{work}
  Q_1 &= \dfrac{4 + 4}{2} = 4 \qquad Q_3 = \dfrac{11 + 13}{2} = 12
\end{work}

(d) Five-number summary: \ans{$2$}, \ans{$4$}, \ans{$7$}, \ans{$12$},
\ans{$16$} \quad IQR: \ans{$8$ games}

% ---------------------------------------------------------------- C4 (2.2)
\boxguard[22]
\item The Harlow school store sold these numbers of water bottles on five days:
$15$, $20$, $21$, $21$, $23$. The mean is $20$ bottles.

\vspace{2pt}
(a) Complete the table.

\vspace{2pt}
\renewcommand{\arraystretch}{1.3}
\begin{tabularx}{\linewidth}{@{} >{\bfseries}l *{5}{>{\centering\arraybackslash}X} | c @{}}
\rowcolor{frost}
 & & & & & & \textbf{Sum} \\
Value $x$ & $15$ & $20$ & $21$ & $21$ & $23$ & $100$ \\
Deviation $x - \bar{x}$ & \ans{$-5$} & \ans{$0$} & \ans{$1$} & \ans{$1$} &
  \ans{$3$} & \ans{$0$} \\
Squared $(x - \bar{x})^2$ & \ans{$25$} & \ans{$0$} & \ans{$1$} & \ans{$1$} &
  \ans{$9$} & \ans{$36$} \\
\end{tabularx}

\vspace{2pt}
(b) Find the variance and the standard deviation.

\begin{work}
  s^2 &= \dfrac{36}{5 - 1} = 9 \text{ square bottles} \qquad s = \sqrt{9} = 3 \text{ bottles}
\end{work}

(c) Interpret $s$ in context.

\par\ansline{A typical day's sales are about 3 bottles from the 20-bottle mean.}

% ---------------------------------------------------------------- C5 (2.3)
\boxguard[20]
\item The five-number summary of $30$ Harlow students' text messages sent on Monday is
$14$, $40$, $51$, $60$, $98$. The three most extreme values are $14$, $86$, and $98$.

\vspace{2pt}
(a) Find the IQR and both fences.

\begin{work}
  \text{IQR} &= 60 - 40 = 20 \\
  \text{lower fence} &= 40 - 1.5(20) = 10 \qquad \text{upper fence} = 60 + 1.5(20) = 90
\end{work}

(b) Which of $14$, $86$, and $98$ are outliers? \ans{Only $98$}

\vspace{3pt}
(c) In a modified boxplot, where does the upper whisker end? \ans{At $86$}

\vspace{3pt}
(d) The value $14$ looks far from the box. Explain whether it is an outlier.

\par\ansline{No --- 14 is above the lower fence of 10, so it is inside the fences.}

% ---------------------------------------------------------------- C6 (2.3)
\boxguard[26]
\item The graph shows how many push-ups $100$ Harlow students did without stopping.

\vspace{2pt}
\begin{minipage}[t]{0.50\linewidth}
\vspace{0pt}
\begin{tikzpicture}
\begin{axis}[width=7.4cm, height=4.8cm, xmin=0, xmax=30, ymin=0, ymax=100,
  xtick={0,5,...,30}, ytick={0,25,50,75,100}, yticklabel={\pgfmathprintnumber{\tick}\%},
  grid=major, grid style={linegray}, tick label style={font=\scriptsize},
  xlabel={Push-ups without stopping}, ylabel={Cumulative \%}, label style={font=\scriptsize}]
  \addplot[cerulean, thick, mark=*, mark size=1.3pt]
    coordinates {(0,0) (5,6) (10,25) (15,50) (20,75) (25,90) (30,100)};
\end{axis}
\end{tikzpicture}
\end{minipage}\hfill
\begin{minipage}[t]{0.47\linewidth}
\vspace{2pt}
(a) $Q_1$: \ans{$10$} \quad Median: \ans{$15$}

\vspace{4pt}
$Q_3$: \ans{$20$} \quad IQR: \ans{$10$}

\vspace{4pt}
(b) A student did $25$ push-ups. What percentile is that? \ans{$90$th}

\vspace{4pt}
(c) Interpret your answer to (b) in context.
\end{minipage}

\par\ansline{90\% of these 100 students did 25 or fewer push-ups without stopping.}

% ---------------------------------------------------------------- C8 (2.4)
\boxguard[26]
\item Harlow surveyed freshmen and sophomores about hours of homework in one week. The
parallel boxplots share one scale.

\vspace{2pt}
\begin{center}
\begin{tikzpicture}[x=1.05cm, y=1cm]
  \draw[linegray] (1.5,0) -- (11,0);
  \foreach \t in {2,3,...,11} {
    \draw[linegray] (\t,-0.08) -- (\t,0.08);
    \node[below, font=\scriptsize] at (\t,-0.08) {\t};}
  \foreach \t in {2.5,3.5,...,10.5} \draw[linegray] (\t,-0.05) -- (\t,0.05);
  \node[below, font=\scriptsize] at (6.5,-0.45) {Hours of homework};
  \drawbox{0.7}{2}{4}{5}{7}{10}{Freshmen}
  \drawbox{1.6}{3}{5.5}{6}{6.5}{8}{Sophomores}
\end{tikzpicture}
\end{center}

(a) Complete the table.

\vspace{2pt}
\renewcommand{\arraystretch}{1.3}
\begin{tabularx}{\linewidth}{@{} X c c @{}}
\rowcolor{frost}
 & \textbf{Median} & \textbf{IQR} \\
Freshmen & \ans{$5$ h} & \ans{$3$ h} \\
Sophomores & \ans{$6$ h} & \ans{$1$ h} \\
\end{tabularx}

\vspace{4pt}
(b) Compare the two groups' homework, with a comparative word for center \emph{and} for spread.

\par\ansline{Sophomores' median (6 h) is greater than freshmen's (5 h), and freshmen's}
\par\ansline{homework is more spread out (IQR 3 h vs.\ 1 h), so sophomores are more consistent.}

\end{enumerate}
\normalsize

% ======================================================================
\boxguard[10]
\parthead{Part D --- Extended Response \hfill 3 questions, 10 points each}

\small
\begin{enumerate}[leftmargin=2.2em, itemsep=10pt, topsep=2pt, start=23]

% ---------------------------------------------------------------- D1 (2.1, 2.3)
\item Sixty Harlow students reported how many minutes they spent on social media last
Saturday. The histogram and the calculator's 1-Var Stats are below.

\vspace{2pt}
\begin{minipage}[t]{0.56\linewidth}
\vspace{0pt}
\begin{tikzpicture}
\begin{axis}[width=8.2cm, height=4.6cm, ybar interval, xmin=0, xmax=420, ymin=0, ymax=20,
  xtick={0,60,...,420}, x tick label as interval=false, ytick={0,5,...,20}, ymajorgrids, grid style={linegray},
  tick label style={font=\scriptsize},
  scaled x ticks=false, xlabel={Minutes on social media}, ylabel={Students},
  label style={font=\scriptsize}]
  \addplot[fill=frost, draw=cerulean] coordinates
    {(0,17) (60,15) (120,11) (180,8) (240,4) (300,3) (360,2) (420,2)};
\end{axis}
\end{tikzpicture}
\end{minipage}\hfill
\begin{minipage}[t]{0.40\linewidth}
\vspace{4pt}
\renewcommand{\arraystretch}{1.15}
\begin{tabular}{@{} l r @{}}
\rowcolor{frost}
\multicolumn{2}{@{}l@{}}{\textbf{1-Var Stats}} \\
$\bar{x}$ & $134$ \\
Sx & $99$ \\
$n$ & $60$ \\
minX & $5$ \\
$Q_1$ & $52$ \\
Med & $108$ \\
$Q_3$ & $196$ \\
maxX & $405$ \\
\end{tabular}
\end{minipage}

\vspace{4pt}
(a) Name the shape of the distribution, and justify it from the histogram.

\par\ansline{Skewed right: most students spent under 2 hours, and the long tail}
\par\ansline{stretches toward the larger times, out to about 400 minutes.}

(b) The mean ($134$ minutes) is larger than the median ($108$ minutes). Explain why, using the shape.

\par\ansline{The mean uses every value, so the few students in the long right tail pull}
\par\ansline{it up; the median uses only the middle position, so it stays lower.}

\boxguard[7]
(c) Which center and which spread should the school report? Give their values in context and
say why.

\par\ansline{The median, 108 minutes, and the IQR, 144 minutes (the middle half spent 52 to 196).}
\par\ansline{The data are skewed right, so the tail pulls the mean and SD up.}

(d) A student wrote: \emph{``Skewed left --- the tallest bars are on the left, where most
students are.''} Explain what is wrong.

\par\ansline{The skew is named by the long tail, not the peak. The tall bars are on the left,}
\par\ansline{but the tail stretches toward longer times, so it is skewed right.}

% ---------------------------------------------------------------- D2 (2.2)
\boxguard[30]
\item Two Harlow basketball players' points in five games. \quad
\textbf{Ana:} $22$, $23$, $24$, $25$, $26$ ($\bar{x} = 24$). \quad
\textbf{Ben:} $4$, $10$, $14$, $20$, $27$ ($\bar{x} = 15$, $s \approx 8.9$).

\vspace{3pt}
(a) Complete the table for Ana, then find her standard deviation.

\vspace{2pt}
\renewcommand{\arraystretch}{1.3}
\begin{tabularx}{\linewidth}{@{} >{\bfseries}l *{5}{>{\centering\arraybackslash}X} | c @{}}
\rowcolor{frost}
 & & & & & & \textbf{Sum} \\
Points $x$ & $22$ & $23$ & $24$ & $25$ & $26$ & $120$ \\
Deviation $x - \bar{x}$ & \ans{$-2$} & \ans{$-1$} & \ans{$0$} & \ans{$1$} &
  \ans{$2$} & \ans{$0$} \\
Squared $(x - \bar{x})^2$ & \ans{$4$} & \ans{$1$} & \ans{$0$} & \ans{$1$} &
  \ans{$4$} & \ans{$10$} \\
\end{tabularx}

\begin{work}
  s^2 &= \dfrac{10}{5 - 1} = 2.5 \qquad s = \sqrt{2.5} \approx 1.6 \text{ points}
\end{work}

(b) A friend says, \emph{``Ana's standard deviation should be the bigger one --- she scores
way more points.''} Use the formula to explain why the friend is wrong.

\par\ansline{$s$ is built from the deviations --- each game's distance from its own mean.}
\par\ansline{Ana's games sit within 2 points of 24; how big the scores are never enters.}

(c) In the next five games Ana scores exactly $3$ points more than in each of these.
New mean: \ans{$27$} \quad New standard deviation: \ans{$1.6$}

\vspace{3pt}
Explain why the standard deviation came out that way.

\par\ansline{Adding 3 to every game adds 3 to the mean, so not one deviation changes.}

(d) Which player is more consistent? Justify with numbers.

\par\ansline{Ana: $s \approx 1.6$ points against Ben's 8.9, so her games stay close together.}

% ---------------------------------------------------------------- D3 (2.4)
\boxguard[30]
\item Harlow's art department tested $12$ cans of each of two brands of paint. Here are the
drying times, in minutes, from the calculator.

\vspace{3pt}
\renewcommand{\arraystretch}{1.25}
\begin{tabularx}{\linewidth}{@{} X *{7}{c} @{}}
\rowcolor{frost}
\textbf{Brand} & $\bar{x}$ & \textbf{Sx} & \textbf{minX} & $Q_1$ & \textbf{Med} & $Q_3$ &
  \textbf{maxX} \\
\midrule
Brand C & $60$ & $2.4$ & $55$ & $58$ & $60$ & $62$ & $65$ \\
Brand D & $60$ & $8.1$ & $44$ & $54$ & $60$ & $66$ & $76$ \\
\end{tabularx}

\vspace{4pt}
Both distributions are roughly symmetric.

\vspace{3pt}
(a) Check Brand D for outliers with the $1.5 \times \text{IQR}$ rule.

\begin{work}
  \text{IQR} &= 66 - 54 = 12 \\
  \text{fences} &= 54 - 18 = 36 \quad\text{and}\quad 66 + 18 = 84
\end{work}

\par\ansline{Every value lies between the fences (44 to 76), so Brand D has no outliers.}

(b) The department head says, \emph{``Same median, same mean --- it doesn't matter which we
buy.''} Respond.

\par\ansline{When the centers match, the spread decides. Brand D's drying times are much}
\par\ansline{more spread out, so some D cans take far longer than any C can.}

(c) Write a comparison of the two brands --- shape, outliers, center, and spread --- using
comparative words.

\par\ansline{Both brands are roughly symmetric with no outliers, and both centers are 60 minutes.}
\par\ansline{Brand D is much more spread out (Sx 8.1 min, range 32 min) than Brand C}
\par\ansline{(Sx 2.4 min, range 10 min), so Brand C's drying times are more consistent.}

(d) A mural must be dry within $65$ minutes. Which brand should the department buy?
Justify with the table.

\par\ansline{Brand C. Its maximum is 65 minutes, so every C can dried within 65 minutes;}
\par\ansline{Brand D's $Q_3$ is 66, so at least a quarter of D cans take too long.}

\end{enumerate}
\normalsize

\end{document}
